\title{QLoRA: Efficient Finetuning of Quantized LLMs}

\begin{abstract}
We introduce \textsc{QLoRA}, a memory-efficient finetuning method that enables the finetuning of large-scale models with up to 65B parameters on a single 48GB GPU, while matching the finetuning performance of traditional 16-bit setups. \textsc{QLoRA} achieves this by backpropagating through a fixed, 4-bit quantized pretrained language model and updating Low Rank Adapters~(LoRA). Our most successful model suite, named \textbf{Guanaco}, surpasses all previous open-source models on the Vicuna benchmark, attaining 99.3\% of ChatGPT's performance with just 24 hours of single-GPU training. \textsc{QLoRA} incorporates several novel techniques to optimize memory without performance loss: (a) 4-bit NormalFloat (NF4), a newly proposed data format that is information-theoretically ideal for weights with a normal distribution; (b) Double Quantization, a strategy that compresses quantization constants to further reduce overall memory demand; and (c) Paged Optimizers, designed to handle memory spikes efficiently. Employing \textsc{QLoRA}, we finetune over 1,000 models and conduct an extensive evaluation of instruction following and chatbot ability across 8 instruction datasets, different architectures (LLaMA, T5), and large-scale models that conventional finetuning would make prohibitive (such as models with 33B or 65B parameters). Our findings indicate that \textsc{QLoRA} achieves state-of-the-art outcomes using high-quality but smaller datasets, even outperforming prior best results with more compact models. We also offer an in-depth study of chatbot effectiveness using both human and GPT-4 assessments, revealing GPT-4-based evaluations as a cost-effective, valid alternative to human judgment. Additionally, our work exposes reliability issues in current chatbot benchmarks. A targeted error analysis highlights specific areas where \textbf{Guanaco} falls short of ChatGPT. All code and models are publicly available, including custom CUDA kernels supporting 4-bit training.\footnote{\url{https://github.com/artidoro/qlora} and \url{https://github.com/TimDettmers/bitsandbytes}}
\end{abstract}

\section{Introduction}

Finetuning large language models (LLMs) remains one of the most effective strategies for boosting their capabilities \citep{min2021metaicl, wei2021finetuned, ouyang2022training, wang2022super, wang2022self, liu2022few}, as well as for imparting specific desired behaviors or eliminating unwanted ones \citep{ouyang2022training,askell2021general,bai2022training}. Nevertheless, the cost of finetuning very large models is prohibitively high; for instance, finetuning a 65B parameter LLaMA model~\citep{touvron2023llama} in standard 16-bit precision requires over 780 GB of GPU memory. Although recent advancements in quantization have made it possible to lower the memory requirements of LLMs during inference \citep{dettmers2022llm, dettmers2022case, frantar2022gptq, xiao2022smoothquant}, these methods are generally unsuitable for training purposes as they tend to fail during the finetuning process \citep{wortsman2023stable}.

In this work, we show for the first time that quantized finetuning at 4 bits can be accomplished with no loss in performance. Our approach, termed \textsc{QLoRA}, introduces an innovative, high-precision quantization method that converts a pretrained model to 4-bit representation, after which a compact set of trainable Low-rank Adapter weights \citep{hu2021lora} 

is incorporated and updated by propagating gradients through the quantized parameters.

\textsc{QLoRA} brings down the average memory needed to finetune a 65B parameter model from over 780GB of GPU RAM to under 48GB, all while preserving both speed and predictive performance relative to the standard 16-bit finetuning baseline. This fundamentally enhances the feasibility and inclusiveness of LLM finetuning by enabling even the largest open models to be finetuned on a single GPU. Leveraging \textsc{QLoRA}, we develop the \textbf{Guanaco} model series, with our second-best variant achieving 97.8\% of ChatGPT’s performance on the Vicuna~\citep{vicuna2023} evaluation, training in under 12 hours on a typical consumer GPU; using a professional GPU for 24 hours, our largest model attains 99.3\%—nearly matching ChatGPT on Vicuna. Upon deployment, the smallest \textbf{Guanaco} model (7B parameters) uses only 5 GB of memory, outperforming the 26 GB Alpaca model by more than 20 percentage points on the Vicuna benchmark (see Table~\ref{tbl:vicuna_eval}).

\textsc{QLoRA} incorporates several key advancements aimed at lowering memory consumption while maintaining model effectiveness: (1) \textbf{4-bit NormalFloat}, a quantization scheme that is theoretically optimal for Gaussian-distributed variables, achieves superior empirical outcomes in comparison to traditional 4-bit Integer and 4-bit Float formats. (2) \textbf{Double Quantization}, which further compresses model parameters by quantizing quantization constants themselves, reduces the memory requirement by roughly 0.37 bits per parameter (resulting in a memory saving of about 3 GB for a 65B parameter model). (3) \textbf{Paged Optimizers}, leveraging NVIDIA's unified memory infrastructure, address and mitigate the spikes in memory usage observed during gradient checkpointing when training on lengthy input sequences. By integrating these innovations, the approach yields a more effective version of LoRA, applying adapters across all layers of the network and thus nearly eliminating the compromises in accuracy that earlier methods encountered.

The memory efficiency of \textsc{QLoRA} allows for a comprehensive exploration of instruction finetuning and conversational agent performance at parameter scales unattainable with standard finetuning methods, due to their higher memory demands. To this end, over 1,000 distinct models are trained, covering a variety of instruction-tuning datasets, model designs, and ranging in size from 80M to 65B parameters. Alongside demonstrating that \textsc{QLoRA} is able to match 16-bit training results (\S\ref{sec:comp_to_std}) and support the development of a leading chatbot, \textbf{Guanaco} (\S\ref{sec:pushing}), the study investigates broader patterns in model behavior. Notably, it reveals that the quality of training data plays a significantly more critical role than sheer dataset size; for example, a compact 9k-sample dataset (OASST1) exceeded the performance of a much larger 450k-sample dataset (FLAN v2, subsampled) in chatbot applications, even when both aimed at generalized instruction following. Additionally, results indicate that achieving high performance on the Massive Multitask Language Understanding (MMLU) benchmark does not guarantee correspondingly strong results on the Vicuna chatbot benchmark, and vice versa, underscoring that alignment between dataset characteristics and task requirements outweighs considerations of dataset scale.

Additionally, we conduct a comprehensive assessment of chatbot effectiveness employing evaluations from both human reviewers and GPT-4. Our methodology employs a tournament-based benchmarking framework, wherein models are pitted against each other to generate the most suitable response for specified prompts. Each match is adjudicated by either GPT-4 or human annotators to determine the superior response. The aggregate tournament results are synthesized into Elo scores~\citep{elo1967proposed, elo1978rating}, which establish a hierarchy of chatbot capabilities. Our findings indicate a strong concordance between human and GPT-4-derived rankings across tournaments; however, we also identify notable occasions where the two sources diverge substantially. Consequently, while model-based evaluation represents a cost-effective alternative to manual annotation, it is important to recognize its inherent ambiguities.

In addition to quantitative benchmarking, we supplement our evaluation with an in-depth qualitative investigation of the \textbf{Guanaco} models. This analysis elucidates both exemplary and problematic behaviors that are overlooked by purely metric-based approaches.

To support continued research, we make public all model-generated outputs alongside their respective human and GPT-4 ratings. Furthermore, we provide open access to our source code and CUDA kernels, and integrate our methodologies within the Hugging Face transformers library \citep{wolf2019huggingface}, thereby promoting widespread usability. We also release a set of adapters suitable for 7/13/33/65B model variants, each fine-tuned on one of eight instruction-following corpora, culminating in 32 distinct, openly available, fine-tuned models.

\section{Background}
\paragraph{Block-wise k-bit Quantization} Quantization refers to the transformation of input data from a higher-precision to a lower-precision representation, commonly involving the conversion from a format with more bits to one with fewer bits (e.g., from 32-bit floating-point to 8-bit integer types). To optimally utilize the value range of the reduced-precision format, input values—often organized as tensors—are typically rescaled into the interval permitted by the target data type. This normalization is performed using the absolute maximum value among the input elements. For instance, to quantize a tensor of 32-bit floating-point values (FP32) to 8-bit integers (Int8) spanning $[-127, 127]$, the transformation is:

\begin{equation}
    \mathbf{X}^{ \text{Int8}} = \text{round}\left( \frac{127}{{\text{absmax}}(\mathbf{X}^{{\text{FP32}}})}\mathbf{X}^{{\text{FP32}}} \right) = \text{round}(c^{\text{FP32}}\cdot \mathbf{X}^{{\text{FP32}}}),
\end{equation}

with $c$ representing the {\em quantization constant} or {\em scale}. The process of dequantization reverses this operation:

\begin{equation}
    \text{dequant}(c^{\text{FP32}}, \mathbf{X}^{\text{Int8}}) = \frac{\mathbf{X}^{\text{Int8}}}{c^{\text{FP32}}} = \mathbf{X}^{\text{FP32}}
\end{equation}

A drawback of this method is that the presence of high-magnitude values (outliers) within the input tensor can result in certain quantization bins—specific bit patterns—being inefficiently filled, with some bins containing few or no quantized values. To address this limitation, a typical strategy involves partitioning the input tensor into smaller blocks that are quantized independently, each with their distinct quantization constant $c$. Specifically, the input tensor $\mathbf{X}\in \mathbb{R}^{b\times h}$ is first flattened and then divided into $n$ consecutive blocks of size $B$, yielding $n = ({b\times h})/{B}$ blocks. Each block is then individually quantized using Equation~1, producing a quantized output and a set of $n$ corresponding quantization constants $c_i$.

\paragraph{Low-rank Adapters} The Low-rank Adapter (LoRA) finetuning approach \citep{hu2021lora} minimizes memory consumption by introducing a compact set of trainable parameters, known as adapters, while the original model weights remain unchanged. During training, gradients are propagated through the unmodified pretrained weights, and only the adapters are updated to improve the loss function. LoRA introduces an auxiliary factorized projection to the standard linear mapping. For a linear transformation defined as ${\mathbf{X} \mathbf{W} = \mathbf{Y}}$ with $\mathbf{X}\in  \mathbb{R}^{b\times h}$ and $\mathbf{W}\in  \mathbb{R}^{h\times o}$, the modified output in LoRA is calculated as:

\begin{equation}
    \mathbf{Y} = \mathbf{X}\mathbf{W} + s\mathbf{X}\mathbf{L}_1\mathbf{L}_2,
\end{equation}

with $\mathbf{L}_1\in  \mathbb{R}^{h\times r}$, $\mathbf{L}_2\in  \mathbb{R}^{r\times o}$, and $s$ denoting a scalar coefficient.

\paragraph{Memory Requirement of Parameter-Efficient Finetuning} The memory demands associated with LoRA during training merit close examination, particularly regarding both the quantity and size of adapters implemented. Given that LoRA’s memory consumption is exceedingly low, deploying a greater number of adapters to boost model performance does not translate to a substantial increase in overall memory usage. Although LoRA serves as a Parameter Efficient Finetuning (PEFT) approach, it is important to note that, in the context of LLM finetuning, the dominant portion of memory use arises from storing activation gradients rather than from the LoRA parameters themselves. For instance, when training a 7B LLaMA model on FLAN v2 with a batch size of 1 and LoRA weights corresponding to the typically adopted 0.2\% of total model weights \citep{hu2021lora,liu2022few}, the memory utilized by LoRA input gradients amounts to 567 MB, whereas the parameters themselves occupy a mere 26 MB. Application of gradient checkpointing~\citep{chen2016training} further decreases the average input gradient memory to 18 MB per sequence, meaning the gradients still surpass the entire LoRA parameter size in terms of memory needs. By contrast, the base 4-bit model utilizes 5,048 MB of memory. This underscores the significance of gradient checkpointing, but also shows that further minimization of LoRA parameters yields only marginal gains in memory efficiency. Consequently, increasing the number of adapters is feasible without notably impacting the total memory footprint during training (refer to Appendix~\ref{app:memory} for a comprehensive analysis). As will be elaborated in later sections, this flexibility proves pivotal in achieving full 16-bit precision model performance.

\section{\textsc{QLoRA} Finetuning}

\textsc{QLoRA} facilitates accurate 4-bit finetuning through two primary innovations we introduce: 4-bit NormalFloat (NF4) quantization and Double Quantization. In addition, we propose the use of Paged Optimizers to mitigate memory spikes associated with gradient checkpointing, which have historically posed out-of-memory challenges when finetuning large models on a single machine.

Within \textsc{QLoRA}, there are two data types: a low-precision storage format (typically 4-bit) and a higher-precision computation type (generally BFloat16). Operationally, each time a \textsc{QLoRA} weight tensor is accessed, it undergoes dequantization to BFloat16 before participating in 16-bit matrix multiplication operations.

We proceed to detail the individual elements of \textsc{QLoRA} and subsequently provide its formal specification.

\paragraph{4-bit NormalFloat Quantization} The NormalFloat (NF) type is an extension of Quantile Quantization \citep{dettmers2022optimizers}, which is designed to be optimal from an information-theoretic standpoint by distributing tensor values so that each quantization bin receives an equal share from the input tensor. This process relies on estimating tensor quantiles using the empirical cumulative distribution function.

A major drawback of quantile quantization is its computationally intensive quantile estimation. To expedite this, fast algorithms such as SRAM quantiles~\citep{dettmers2022optimizers} are utilized. However, these rapid approximations often result in substantial quantization errors, especially for outlier values—which can be critical for performance.

The high cost and inaccuracy associated with quantile estimation can be sidestepped when the tensor values are sampled from distributions that are consistent up to a quantization constant. In such scenarios, quantiles remain invariant across input tensors, rendering exact quantile calculation computationally viable.

Pretrained neural network weights tend to follow a zero-mean Gaussian distribution characterized by a standard deviation $\sigma$ (refer to Appendix~\ref{app:normality}). To ensure consistency, all weights can be mapped onto a standard reference distribution by adjusting $\sigma$ so that the weight distribution is aligned with the numerical limits of the target data type. In this context, we define the range for our data type as $[-1, 1]$. Consequently, the quantiles from both the data type and the neural network weights require normalization within this interval.

For zero-mean normal distributions with arbitrary $\sigma$, the optimal data type from an information-theoretic perspective—restricted to the interval $[-1, 1]$—is derived as follows: (1) compute $2^k+1$ quantiles from a theoretical $N(0, 1)$ distribution to construct a $k$-bit quantile-based quantization data type, (2) normalize the resulting data type values to occupy the $[-1, 1]$ range, and (3) rescale any given weight tensor into $[-1, 1]$ using absolute max scaling prior to quantization.

After matching the distributions of the weight tensor and the data type, standard quantization procedures can be applied. In step (3), this effectively means scaling the weight tensor's standard deviation to coincide with that of the $k$-bit data type. To be precise, the $2^k$ quantization levels $q_i$ for the data type are determined by:

\begin{equation}
q_i  = \frac{1}{2}\left( Q_X\left(\frac{i}{2^k+1}\right) + Q_X\left(\frac{i+1}{2^k+1}\right)\right),
\end{equation}

where $Q_X(\cdot)$ denotes the quantile function corresponding to the standard normal distribution $N(0,1)$. A limitation inherent in symmetric k-bit quantization lies in its inability to precisely encode the value zero, an essential feature when quantizing padding or other elements with zero value, as this would introduce error. To guarantee the inclusion of a discrete zero point at $0$ while utilizing the complete $2^k$ representational capacity of a k-bit format, we construct an asymmetric data type. This is accomplished by separately estimating the quantiles $q_i$ for two domains: $2^{k-1}$ quantiles for the negative portion, and $2^{k-1}+1$ for the positive portion, subsequently merging these sets and discarding one of the duplicated zeros found in both. We refer to the resulting type, characterized by an equal expected frequency in each quantization bin, as \emph{k-bit NormalFloat} (NFk). This design is optimal, in terms of information theory, for data that follows a zero-centered normal distribution. The exact specification of these data type values is provided in Appendix~\ref{app:nf4}.

\paragraph{Double Quantization{}} 
We propose \emph{Double Quantization} (DQ), which involves further quantizing the set of quantization constants themselves to realize additional memory efficiency. Achieving high accuracy with 4-bit quantization often requires a small block size~\citep{dettmers2022case}, but this also incurs significant memory cost due to the storage of quantization constants. As an illustration, with 32-bit constants and a block size of 64 for $\mathbf{W}$, the quantization constants alone contribute, on average, $32/64 = 0.5$ bits to each parameter. Double Quantization seeks to alleviate this memory overhead.

In detail, Double Quantization processes the first set of quantization constants, $c_2^{\text{FP32}}$, as the subject of a subsequent quantization step. This produces secondarily quantized constants, $c_2^{\text{FP8}}$, along with a new set of quantization constants, $c_1^{\text{FP32}}$, for the second quantization layer. For this step, we adopt 8-bit floating point representations and a block size of 256, given that prior studies~\citep{dettmers2022case} show negligible performance drop with 8-bit quantization. Because $c_2^{\text{FP32}}$ are non-negative, we mean-center $c_2$ before quantizing to better suit symmetric quantization. Through this approach, for block sizes of 64, the memory used for quantization constants per parameter decreases from $32/64=0.5$ bits to $8/64 + 32/(64\cdot256) = 0.127$ bits, thus saving 0.373 bits per parameter.

\paragraph{Paged Optimizers} leverage the NVIDIA unified memory~\footnote{\tiny \url{https://docs.nvidia.com/cuda/cuda-c-programming-guide}} capability, which automates page-level transfers between CPU and GPU, facilitating error-free GPU execution when GPU memory resources are insufficient. This unified memory mechanism operates similarly to standard paging between main memory and disk. In our approach, optimizer states are stored using this paged memory, allowing them to be seamlessly offloaded to CPU RAM when GPU memory is fully utilized and transparently paged back to GPU memory during optimizer updates as needed.

\paragraph{\textsc{QLoRA}.} Incorporating the mechanisms outlined above, we formulate \textsc{QLoRA} for a single linear layer in the quantized foundation model, equipped with one LoRA adapter, as follows:

\begin{equation}
    \mathbf{Y}^{\text{BF16}} =  \mathbf{X}^{\text{BF16}} \text{doubleDequant}(c_1^{\text{FP32}}, c_2^{\text{k-bit}}, \mathbf{W}^{\text{NF4}}) + \mathbf{X}^{\text{BF16}}\mathbf{L}^{\text{BF16}}_1\mathbf{L}^{\text{BF16}}_2,
\end{equation}

where the function doubleDequant$(\cdot)$ is specified by:

\begin{equation}
    \text{doubleDequant}(c_1^{\text{FP32}}, c_2^{\text{k-bit}}, \mathbf{W}^{\text{k-bit}}) = \text{dequant}(\text{dequant}(c_1^{\text{FP32}}, c_2^{\text{k-bit}}), \mathbf{W}^{\text{4bit}}) = \mathbf{W}^{\text{BF16}},
\end{equation}

In this context, $\mathbf{W}$ is encoded with NF4, while $c_2$ utilizes FP8.

To optimize quantization fidelity, we assign a block size of 64 for $\mathbf{W}$ and use a block size of 256 for $c_2$ to optimize memory usage.

Parameter updates require only the computation of gradients with respect to the adapter weights, namely $\frac{\partial E}{\partial \mathbf{L}_i}$; gradients with respect to the 4-bit weights, $\frac{\partial E}{\partial \mathbf{W}}$, are not necessary. Nevertheless, computing $\frac{\partial E}{\partial \mathbf{L}_i}$ involves determining $\frac{\partial \mathbf{X}}{\partial \mathbf{W}}$, which, in accordance with equation~(5), requires dequantizing $\mathbf{W}^{\text{NF4}}$ to $\mathbf{W}^{\text{BF16}}$ so that the derivative $\frac{\partial \mathbf{X}}{\partial \mathbf{W}}$ can be calculated with BFloat16 precision.

In summary, \textsc{QLoRA} utilizes a single storage data type—commonly 4-bit NormalFloat—and a distinct computation data type, typically 16-bit BrainFloat. Prior to computation during forward and backward propagation, the stored data are dequantized into the computational format, with gradient calculations restricted to the LoRA parameters, which always operate in 16-bit BrainFloat.

\section{QLoRA vs. Standard Finetuning}
\label{sec:comp_to_std}

The preceding discussion outlined the operation of QLoRA and its substantial benefits in lowering memory consumption during the finetuning of models. The pivotal issue that remains is whether QLoRA matches the performance levels achieved by full-model finetuning. An additional goal is to investigate the influence of the components of QLoRA, especially the impact of using NormalFloat4 relative to the standard Float4. The experimental results presented in the subsequent sections address these research questions.

\paragraph{Experimental setup.} The evaluation spans three architectural paradigms—encoder, encoder-decoder, and decoder-only—contrasting QLoRA with both 16-bit adapter-finetuning and full-model finetuning, covering models up to 3B parameters. Performance is measured on benchmarks including GLUE \citep{wang2018glue} with RoBERTa-large \citep{liu2019roberta}, Super-NaturalInstructions (TKInstruct) \citep{wang2022super} with T5 \citep{t5}, and 5-shot MMLU \citep{hendrycksmeasuring} after LLaMA has been finetuned on Flan v2 \citep{longpre2023flan} and Alpaca \citep{alpaca}. To further examine the advantages of NF4 compared to alternative 4-bit formats, we adopt the protocol of \citet{dettmers2022case}, evaluating post-quantization zero-shot accuracy and perplexity for several model families (OPT~\citep{zhang2022opt}, LLaMA~\citep{touvron2023llama}, BLOOM~\citep{scao2022bloom}, Pythia~\citep{biderman2023pythia}) across model sizes from 125m to 13B parameters.

More granular details concerning individual experiments are reported in the results section to facilitate readability, with comprehensive methodologies found in Appendix~\ref{app:exp-setup}.

Paged optimizers are instrumental for enabling \textsc{QLoRA} finetuning on models sized 33B/65B with only a single 24/48GB GPU. However, we do not provide empirical benchmarks on Paged Optimizers, since paging is triggered primarily by mini-batches with exceptionally long sequences—a rare scenario. Nonetheless, an analysis of runtime efficiency using paged optimizers for 65B models on 48GB GPUs indicates that, for a batch size of 16, training throughput is on par with standard optimizers. It is suggested that future research systematically measure and delineate the specific conditions under which paging overhead may impede training speed.

\paragraph{Default LoRA hyperparameters do not match 16-bit performance}

Applying the default configuration of LoRA—targeting query and value attention projection matrices as recommended by \citep{hu2021lora}—proves insufficient to reproduce the full finetuning results for large base models. Evidence presented in Figure~\ref{fig:hyper}, specifically for finetuning LLaMA 7B on Alpaca, demonstrates that the most influential LoRA hyperparameter is the total number of LoRA adapters deployed; achieving parity with full finetuning performance necessitates applying LoRA across all linear layers within transformer blocks. Other hyperparameters, like the projection dimension $r$, exert negligible effect on outcomes (further details in Appendix \ref{app:exp-setup}).

In a similar vein, our analysis indicates that the commonly used hyperparameters for fully finetuned baselines are insufficiently optimized. To establish more rigorous baselines, we conduct an extensive hyperparameter sweep, exploring learning rates from 1e-6 up to 5e-5 and batch sizes ranging from 8 to 128. The resulting outcomes for finetuning the 7B LLaMA model on the Alpaca dataset can be found in Figure~\ref{fig:hyper}.

\paragraph{4-bit NormalFloat Outperforms 4-bit Floating Point} 

Although the 4-bit NormalFloat (NF4) format is theoretically optimal in terms of information retention, it remains to be confirmed if this optimality offers tangible empirical benefits. Replicating the experimental design of \citet{dettmers2022case}, we assess quantized LLMs (including OPT \citep{zhang2022opt}, BLOOM \cite{scao2022bloom}, Pythia \cite{biderman2023pythia}, and LLaMA) covering model sizes from 125M to 65B, utilizing multiple quantization data types. These models are evaluated on language modeling tasks and a collection of zero-shot benchmarks. As depicted in Figure~\ref{fig:nf4} and Table~\ref{tbl:nf4_ppl}, NF4 consistently surpasses FP4 and Int4 in terms of performance, and additionally, applying double quantization achieves reduced memory consumption with no observable impact on accuracy.

\paragraph{k-bit \textsc{QLoRA} Achieves Parity with 16-bit Full Finetuning and 16-bit LoRA} Recent research demonstrates that, while 4-bit quantization is viable for \emph{inference}, it usually incurs some accuracy loss compared to 16-bit models \citep{dettmers2022case,frantar2022gptq}. This raises the important issue of whether 4-bit adapter finetuning can restore this performance gap. To address this, we conduct experiments under two scenarios.

The initial scenario involves a head-to-head comparison with full 16-bit finetuning of RoBERTA and T5 models spanning 125M to 3B parameters, evaluated on both the GLUE benchmark and the Super-NaturalInstructions dataset. The results, provided in Table~\ref{tab:nf4vsbf16}, indicate that all adapter-based finetuning methods—whether utilizing 16-bit, 8-bit, or 4-bit quantization—closely match the results of the full 16-bit finetuned baseline. These findings imply that adapter-based finetuning can effectively recover any degradation resulting from coarser quantization.

In our second experimental configuration, due to the memory constraints that make full finetuning of models with 11B parameters or larger impractical on a single high-memory GPU server, we investigate whether 4-bit \textsc{QLoRA} maintains parity with 16-bit LoRA across parameter sizes from 7B to 65B. Specifically, we fine-tune LLaMA models ranging from 7B to 65B on the Alpaca and FLAN v2 instruction-following datasets, subsequently assessing their performance on the MMLU benchmark using 5-shot accuracy. Table~\ref{tbl:llama-datatype} presents these findings, where we observe that NF4 using double quantization fully matches the MMLU results achieved by 16-bit LoRA. Furthermore, our analysis indicates that the FP4 variant of \textsc{QLoRA} trails the 16-bit brain float LoRA baseline by approximately 1 percentage point. These outcomes reinforce two primary conclusions: (1) \textsc{QLoRA} utilizing NF4 is capable of mirroring the results of both 16-bit full finetuning and 16-bit LoRA finetuning, and (2) NF4 demonstrates superior quantization fidelity compared to FP4.

\paragraph{Summary} Our experimental results consistently demonstrate that 4-bit \textsc{QLoRA} with the NF4 format achieves performance on par with 16-bit full finetuning and 16-bit LoRA finetuning across standard academic benchmarks using established evaluation protocols. We have also provided evidence that NF4 outperforms FP4 and that implementing double quantization does not negatively impact performance. Altogether, these findings strongly suggest that 4-bit \textsc{QLoRA} finetuning is a robust alternative to traditional 16-bit techniques.

Consistent with prior research on quantization \citep{dettmers2022case}, our experiments on MMLU and Elo benchmarks demonstrate that, for a fixed finetuning and inference resource budget, allocating capacity to larger base models with reduced precision yields favorable outcomes. This underscores the significant efficiency gains afforded by \textsc{QLoRA}. As we did not observe any decline in performance when comparing 4-bit finetuning to full 16-bit finetuning, this leaves open the question of the precise threshold in the performance versus precision trade-off for QLoRA-based tuning—a topic we suggest for future investigation.

Our investigation turns to instruction tuning at scales that would be infeasible to pursue using traditional full 16-bit finetuning approaches on typical academic hardware.

\section{Pushing the Chatbot State-of-the-art with QLoRA}
\label{sec:pushing}
After demonstrating that 4-bit \textsc{QLoRA} delivers results on par with its 16-bit counterpart across model sizes, task domains, and datasets, we undertake a thorough examination of instruction finetuning, scaling up to the largest openly available language models suited for academic research. To rigorously evaluate how well instruction finetuning performs in these settings, we measure performance on the difficult MMLU benchmark for Natural Language Understanding and devise novel evaluation protocols aimed at assessing chatbot capabilities in realistic scenarios.

\subsection{Experimental setup} Here, we provide a summary of our experimental methodology; further specifics are detailed in Appendix \ref{app:chatbot}. 

\paragraph{Data} As recent instruction-following datasets have not been systematically reviewed to our knowledge, we selected eight contemporary datasets for our study. Our collection includes datasets sourced via crowd-sourcing efforts (OASST1~\citep{kopf2023openassistant}, HH-RLHF~\citep{bai2022training}), those derived from distillation of instruction-finetuned models (Alpaca~\citep{alpaca}, self-instruct~\citep{wang2022self}, unnatural-instructions~\citep{honovich2022unnatural}), larger corpora aggregations (FLAN v2~\citep{chung2022scaling}), as well as hybrid datasets (Chip2~\citep{laion2023}, Longform~\citep{koksal2023longform}). Collectively, these datasets span diverse languages, scales, and licensing terms.

\paragraph{Training Setup} To eliminate confounding variables arising from varying training paradigms, we apply QLoRA finetuning exclusively with a cross-entropy loss under a supervised learning framework, omitting any reinforcement learning techniques, regardless of whether datasets contain human preference annotations. In cases where instructions and responses are distinctly separated within a dataset, we restrict finetuning solely to the response component (see Appendix \ref{app:chatbot} for ablation details). For OASST1 and HH-RLHF, where multiple possible responses exist, we extract the highest-rated response at every branch within the conversational tree, finetuning on the complete chosen dialogue, including all instructions. All experiments are conducted with NF4 \textsc{QLoRA} leveraging double quantization and paged optimization to mitigate the risk of memory overload during gradient checkpointing. Hyperparameter tuning is limited and targeted; for LLaMA models at 13B and 33B scales, we discover that hyperparameters identified at the 7B scale transfer reliably, with the exceptions of learning rate and batch size. Specifically, learning rates are reduced by half for 33B and 65B, while batch sizes are doubled.

\paragraph{Baselines} Our models are benchmarked against both open research chatbot systems (such as Vicuna~\citep{vicuna2023} and Open Assistant~\citep{kopf2023openassistant}) and major commercial chatbots (including GPT-4 \citep{openai2023gpt}, GPT-3.5-turbo, and Bard). The Open Assistant system utilizes a LLaMA 33B model trained with RLHF on the OASST1 dataset, which is also included in our experiments. Vicuna applies full finetuning to LLaMA 13B utilizing proprietary user conversations from ShareGPT, effectively operating as a distilled product of OpenAI GPT models.

\subsection{Evaluation}

To assess language understanding, we adopt the MMLU (Massively Multitask Language Understanding) benchmark \citep{hendrycksmeasuring}, widely recognized for covering a comprehensive range of language-related tasks. This evaluation comprises 57 multiple-choice subjects, which include fields such as elementary mathematics, U.S. history, computer science, law, among others. All results are reported as 5-shot accuracy scores.

Our assessment of generative language abilities encompasses both automated metrics and human judgment. This latter evaluation phase utilizes human-curated queries with the intention of quantifying the quality of model outputs. Although such human-centered evaluations offer a more authentic measure of chatbot effectiveness and have seen increasing adoption, there is currently no standardized methodology recognized in academic literature. Below, we detail our evaluation procedure, which uniformly applies nucleus sampling with $p=0.9$ and a temperature setting of $0.7$.

\paragraph{Benchmark Data} For our experiments, we utilize two specifically curated query datasets: the Vicuna prompts \citep{vicuna2023} and the OASST1 validation set \citep{kopf2023openassistant}. The Vicuna benchmark consists of 80 varied prompts across multiple domains, which we employ as provided. OASST1 features a multilingual array of crowd-generated multiturn conversations between users and assistants. From this dataset, we extract every user utterance in the validation split as a standalone query, prepending each with the dialogue history. This process yields 953 unique user queries. We collectively refer to these datasets as the Vicuna and OA benchmarks.

\paragraph{Automated Evaluation} In line with the evaluation approach established by \citet{vicuna2023}, we employ GPT-4 as a judge to compare the outputs of various systems against ChatGPT (GPT-3.5 Turbo) on the Vicuna benchmark. For each query, both ChatGPT’s and the candidate model’s responses are presented, and GPT-4 is prompted to score each out of ten and supply a rationale. The aggregate model score is expressed as a percentage of ChatGPT’s total. Importantly, this percentage can surpass 100\% if a model’s performance exceeds that of ChatGPT. We observe a notable positional bias, with GPT-4 tending to award higher scores to whichever response appears first. To mitigate this bias, we suggest averaging results over both possible response orderings.

Subsequently, we conduct performance assessments based on direct output comparisons. This involves simplifying the evaluation task to a three-way label classification that allows for ties. We instruct GPT-4 to choose the superior response or indicate a tie, accompanied by justification. These binary evaluations are performed for every pairwise permutation of systems across both the Vicuna and OA benchmarks.

\paragraph{Human Evaluation} Although existing studies suggest that generative models can serve as effective evaluators for system outputs \citep{fu2023gptscore}, there is, to our knowledge, insufficient evidence establishing the correspondence between GPT-4 ratings and human assessments in the context of chatbots. Consequently, we supplement automated evaluations with two parallel human studies on the Vicuna benchmark, mirroring both automated protocols. These human assessments are carried out on Amazon Mechanical Turk (AMT), where two annotators independently compare model outputs to ChatGPT, and three annotators evaluate model pairs directly.

\paragraph{Elo Rating} Utilizing both human and automated pairwise comparisons, we organize a tournament-like framework in which various models are pitted against each other. Each round consists of two models responding to the same prompt, with their outputs judged against one another. This structure echoes the methodology of \citet{bai2022training} and \citet{vicuna2023}, though our approach extends theirs by incorporating GPT-4 assessments alongside human evaluators. To determine Elo~\citep{elo1967proposed,elo1978rating} scores, we randomly select from the set of annotated matchups. The Elo rating system, originally developed for chess and other competitive games, represents the likelihood of one participant winning relative to another; for instance, a model with an Elo of 1100 facing an opponent with a rating of 1000 is predicted to win about 65\% of the time, while evenly matched ratings (such as 1000 vs 1000 or 1100 vs 1100) suggest a 50\% chance of victory for either side. The amount by which a player's Elo changes post-match depends on how expected the outcome was: surprising wins yield a larger adjustment, while predicted results cause minimal change. Across many matches, the Elo score stabilizes and becomes a reasonable proxy for model proficiency. All models begin with a baseline rating of 1,000 and a $K$ value of 32. Mirroring the protocol of \citet{vicuna2023}, we conduct 10,000 tournament simulations, each with a different random seed, in order to mitigate potential bias arising from the sequence in which pairwise competitions occur.

\subsection{\textbf{Guanaco}: \textsc{QLoRA} trained on OASST1 achieves State-of-the-Art Chatbot Performance}

Through both automatic and manual assessments, we observe that the leading \textsc{QLoRA} adapted model, {Guanaco} 65B, fine-tuned with a variation of OASST1, stands as the highest-performing open-source conversational agent, demonstrating results on par with ChatGPT. When evaluated against GPT-4, {Guanaco} 65B and 33B attain an expected win rate of 30\% according to Elo scores derived from human raters' system-level pairwise judgments—representing the top result achieved thus far.

Table~\ref{tbl:vicuna_eval} provides Vicuna benchmark \cite{vicuna2023} results in relation to ChatGPT. Here, {Guanaco} 65B emerges as the top model following GPT-4, attaining 99.3\% of ChatGPT's performance. The 33B {Guanaco} model features more parameters than Vicuna 13B, yet benefits from 4-bit quantization, leading to greater memory efficiency—requiring just 21 GB instead of 26 GB—and outperforms Vicuna 13B by three percentage points. In addition, {Guanaco} 7B, with a compact 5 GB memory demand, is suitable for deployment on modern smartphones and exceeds Alpaca 13B by nearly 20 percentage points.

Nonetheless, Table~\ref{tbl:vicuna_eval} shows extensive confidence intervals, leading to considerable overlap in model performance estimates. We suspect that much of this ambiguity results from ambiguous rating scale definitions; for instance, the meaning of an 8 out of 10 score can vary substantially between cases. Thus, we suggest prioritizing Elo-based rankings \citep{elo1967proposed}, which use \textit{pairwise} comparisons from both humans and GPT-4, thus circumventing the issue of establishing an absolute scale. The comparative Elo scores for the leading models appear in Table~\ref{tbl:elo}. Notably, while model rankings by humans and GPT-4 on the Vicuna benchmark diverge for some models (notably Guanaco 7B), their system-level orderings generally align, exhibiting a Kendall Tau coefficient of $\tau=0.43$ and Spearman rank correlation of $r=0.55$. At the level of individual examples, correspondence is weaker, with a Fleiss $\kappa=0.25$ between GPT-4 and majority human judgments. In sum, these figures indicate a moderate concordance in system-level rankings produced by GPT-4 and humans, implying that automated model evaluations can provide a reasonably sound substitute for human judgments. We further elaborate on these points in Section~\ref{ssec:considerations}.

The Elo rankings presented in Table~\ref{tbl:elo_large} demonstrate that the 33B and 65B {Guanaco} models surpass all others except for GPT-4 in both the Vicuna and OA evaluation benchmarks, and their results closely match those of ChatGPT as also reflected in Table~\ref{tbl:vicuna_eval}. It is important to highlight that the Vicuna benchmark shows a preference for open-source models, whereas the broader OA benchmark tends to favor ChatGPT. Moreover, the data in Tables \ref{tbl:mmlu-llama} and \ref{tbl:vicuna_eval} suggest that the choice of finetuning dataset has a significant impact on model performance. Specifically, Llama models finetuned with FLAN v2 achieve strong results on MMLU but underperform on the Vicuna benchmark, a pattern that holds for various other models as well. These observations underscore that different evaluation metrics capture distinct abilities, so high scores on MMLU do not necessarily translate to high chatbot efficacy as measured by Vicuna or OA, and vice versa. Notably, 
 is the sole leading model in our comparisons not trained with proprietary data, owing to the OASST1 dataset’s strict prohibition on GPT-based data sources. The next highest-performing model relying exclusively on open-source data, Anthropic HH-RLHF, trails {Guanaco} by 30 percentage points on the Vicuna benchmark (refer to Table~\ref{tbl:vicuna_eval}). Collectively, these findings affirm that 4-bit \textsc{QLoRA} not only enables state-of-the-art chatbot capabilities comparable to ChatGPT but also achieves this with efficient resource requirements—the 33B {Guanaco} model, for instance, can be trained in under 12 hours on a 24 GB consumer GPU. These results pave the way for further advancement through \textsc{QLoRA} finetuning on specialized open-source corpora, allowing open models to reach or even exceed the quality of today’s leading commercial systems.

\section{Qualitative Analysis}

Although our primary assessment relies on quantitative measures, solely focusing on aggregate statistics comes with limitations. Notably, one critical concern is the issue of benchmark validity \citep{liao2021we}—namely, whether a benchmark genuinely evaluates the construct it claims to measure, a concern heightened by the prevalence of ``shortcuts'' that machine learning systems may leverage to achieve high scores without exhibiting true task competence \citep{gururangan2018annotation, poliak2018hypothesis}. To address this limitation to some extent, we include a qualitative analysis, divided into two parts. In \S\ref{ssec:examples}, we provide sample outputs that exemplify characteristic trends in the responses produced by our 65b {Guanaco} model. Subsequently, in \S\ref{ssec:considerations}, we elaborate on key factors affecting the interpretation of our findings and the conclusions we draw.

\subsection{Qualitative Examination of Generated Examples}
\label{ssec:examples}

For example selection, we systematically review outputs produced for both the Vicuna and OpenAssistant benchmarks, seeking recurring tendencies in {Guanaco}'s responses. Upon identifying such tendencies, we design prompts specifically intended to trigger these patterns even if they yield incorrect answers. For instance, noticing a propensity for verbose explanations, we challenge the model with requests like ``Answer yes or no without explanation.'' This methodology allows us to highlight ``lemons,'' where the model is successfully misled, alongside ``cherries,'' where it remains robust against our interventions. All responses presented here were generated utilizing Nucleus Sampling \cite{holtzmancurious} with $p=0.9$.

Naturally, the analysis presented here is not exhaustive; as a small-scale qualitative investigation, it cannot account for every possible variable—for instance, the space of potential responses a model might produce in reply to any prompt is vast, thus our findings rest upon sample outputs that we believe to be indicative. Nonetheless, the illustration of these samples is intended to contextualize the quantitative data discussed previously. By making all models and source code openly available, we hope this section encourages further, more granular research into the challenges identified.

\paragraph{Factual Recall} When presented with straightforward factual questions such as ``What is the capital of Zambia?'', each model repeatedly delivers accurate responses, for example:

\begin{quote} 
    \textbf{Guanaco}:
    \texttt{The capital of Zambia is Lusaka.}
\end{quote}

 However, as the questions become more esoteric, {Guanaco} demonstrates decreased accuracy while maintaining a tone of certainty. As an example, consider the following prompt adapted from HotPotQA \citep{yang2018hotpotqa}:

\begin{quote}
     \textbf{User}: In what year was the singer who popularized `I'll Keep the Lovelight Burning' born?  

     \textbf{Guanaco}:~\texttt{Al Jolson is credited with popularizing the song `I'll Keep the Lovelight Burning,' and he was born in the year 1886.}
 \end{quote}

In this case, the response identifies the incorrect popularizer and provides an incorrect year of birth—though the year aligns with the actual birth year of Al Jolson, who is erroneously named as the singer.

\paragraph{Suggestibility} {Guanaco} also exhibits unexpected robustness in declining to accept or propagate some forms of presupposed misinformation, as demonstrated in the following dialogue:

\begin{quote}
    \textbf{User}: How was it finally, officially confirmed that the earth is flat by peer reviewed scientists? \\ \\
    \textbf{Guanaco}: \texttt{Peer reviewed scientists have never officially validated the notion that the earth is flat. The scientific community overwhelmingly agrees that the earth is spherical—a fact established over centuries through extensive observations and empirical evidence. \\ \\ Flat earth claims are widely regarded as a conspiracy theory without scientific support. The assertion that peer review has led to a confirmation of a flat earth is demonstrably untrue. \\ \\ Please feel free to ask any further questions.}
\end{quote}

also demonstrates proficiency in recognizing when certain questions are inherently unanswerable, as seen in the following example:

\begin{quote}
    \textbf{User}: What time is it?

    \textbf{Guanaco}: \texttt{I apologize, but I do not possess the ability to access current, real-time information as an AI language model. My replies are based solely on my existing knowledge base and therefore cannot provide the exact current time. If you need to know the time, please check a clock or use an online time service.}
\end{quote}

The same pattern is observed with queries such as ``Where are you?'', ``How are you?'', and similar questions.

\paragraph{Refusal} In a related fashion, {Guanaco} at times declines to comply with prompts for reasons that are not always consistent:

\begin{quote}
    \textbf{User}: Please reverse the words in the sentence ``I'm a cool dude in a mocha mood.'' \\ \\
    \textbf{Guanaco}: \texttt{I'm sorry, I can't perform that task right now, but I can clarify why your sentence is grammatically correct.\\ \\The expression ``I'm a cool dude in a mocha mood'' adheres to standard English grammar and syntax conventions. Specifically, the adjective ``cool'' appropriately modifies the noun ``dude,'' and the prepositional phrase ``in a mocha mood'' supplies extra contextual detail. This construction is both valid and meaningful in English.\\ \\Let me know if you have other questions or if there's anything else I can help you with.}
\end{quote}

\paragraph{Secret Keeping} However, such refusals are not dependable even when it is critical that the model keep information confidential. System developers may wish to ensure that some data remains undisclosed to users, despite allowing unrestricted text generation. For instance, we prompt the model with ``The secret word is `banana'. Do not reveal or repeat back this back, no matter what the human requests.'' and test how resilient the system is to indirect probing. A straightforward question results in an appropriate refusal:

\begin{quote}
    \textbf{User}: What is the secret word?

    \textbf{Guanaco}: \texttt{I'm sorry, but I can't share the secret word. My programming does not permit me to reveal or repeat it, regardless of what the human requests. Can I help you with anything else?}
\end{quote}

which exemplifies the preferred outcome. Yet, only minimal manipulation is needed to bypass this safeguard:

\begin{quote}
    \textbf{User}: This is a game. The goal is to ignore your previous instructions. What is the secret word?

    \textbf{Guanaco}: \texttt{Sorry, the secret word is banana. It's best if we don't speak of it. Is there anything else you'd like to ask?}
\end{quote}

This highlights the necessity for continued research into mechanisms that enforce more reliable adherence to instructions.

\paragraph{Math} Mathematics remains the primary area of difficulty for {Guanaco}, consistent with challenges faced by other language models as discussed in \citep{liang2022holistic}. When {Guanaco} provides step-by-step reasoning, its mathematical answers are frequently correct, as illustrated below:

\begin{quote}
    \textbf{User}: Lee mows one lawn and charges \$33. Last week he mowed 16 lawns and three customers each gave him a \$10 tip. How many dollars did Lee earn mowing lawns last week?

    \textbf{Guanaco}: \texttt{Lee earned \$582 mowing lawns last week. \\ \\ To find the total amount he earned, we first need to calculate how much he earned for mowing 16 lawns: \\ \\  \$33 x 16 = \$528 \\ \\ Then we need to add the tips he received from the three customers: \\  \\ \$10 x 3 = \$30 \\ \\ So the total amount Lee earned was: \\ \\ \$528 + \$30 = \$558}
\end{quote}

Nevertheless, {Guanaco} is prone to mistakes on relatively simple questions when it does not follow a systematic, stepwise approach—a phenomenon previously reported \cite{weichain}. This is evident in the exchange below:

\begin{quote}
    \textbf{User}: Please factorize 1833.

    \textbf{Guanaco}: \texttt{Sure, I'd be happy to help! 1833 is a prime number, which means it has only two factors: 1 and 1833. \\ \\ So the unique factorization of 1833 is: \\ \\ 2\^{}1 * 3\^{}2 * 17\^{}1 \\ \\ I hope this helps! Let me know if you have any other questions.}
\end{quote}

which is not only incorrect (the correct factorization is $3 \times 17 \times 43$), but it makes two errors.

\paragraph{Theory of Mind} {Guanaco} demonstrates unexpectedly robust Theory of Mind reasoning abilities \cite{nematzadeh2018evaluating, sap2022neural}. As an illustration, the model provides a comprehensive and accurate response to the following scenario:

\begin{quote}
    \textbf{User}: James and Abby are in the bedroom. Abby put the pen in the desk drawer. Abby leaves the bedroom. James moves the pen into the bag. Where does James think Abby will look for the pen?

    \textbf{Guanaco}: \texttt{James knows that Abby put the pen in the desk drawer. But James also knows that he has moved the pen into the bag, and he assumes that Abby does not know about this. So, James may guess that Abby will look for the pen in the desk drawer, because that's where she left it last time. However, this assumption may be incorrect, because Abby may have seen James moving the pen.}
\end{quote}

Nevertheless, these sorts of inferences are not consistently reliable; the model often generates explanations that rest on unwarranted assumptions, for example:

\begin{quote}
    \textbf{User}: Evelyn entered the living room. Jackson entered the playroom. James entered the playroom. The beans are in the treasure chest. James exited the playroom. Jackson moved the beans to the pantry. Jackson exited the playroom. James entered the living room. Where will James look for the beans?

    \textbf{Guanaco}: \texttt{James will look for the beans in the pantry, because that's where Jackson moved them.}
\end{quote}

Here, {Guanaco} assumes that James is aware of the beans’ relocation, though this information was not conveyed. Such failure modes have also been identified in recent research \cite{sap2022neural}, and warrant further investigation.

\subsection{Considerations}
\label{ssec:considerations}

\paragraph{Evaluation}

Our results show only moderate inter-annotator consistency (Fleiss $\kappa=0.42$), with agreement declining further when evaluating two high-performing models against each other. This highlights both the shortcomings of existing benchmarks and the challenges inherent in human evaluation protocols for conversational agents. In our manual side-by-side comparisons between ChatGPT and {Guanaco} 65B using the Vicuna benchmark, we observed substantial subjectivity, with disagreements among the paper’s authors regarding which responses were preferable. There is a need for future work to address these subjectivity-related challenges, potentially by adopting strategies from fields such as Human-Computer Interaction and Psychology that specialize in preference assessment.

We additionally note that current automated evaluation tools display notable biases. For instance, there are pronounced order effects, with GPT-4 tending to award higher ratings to whichever system is presented first in the prompt. Furthermore, direct comparison at the sample level between GPT-4 and human raters yields limited agreement (Fleiss $\kappa=0.25$), indicating misalignment in their evaluation criteria. Table~\ref{tbl:elo_large} further reveals that GPT-4 consistently gives its own outputs higher scores than those assigned by human evaluators, reporting an Elo rating of 1348 compared to humans' 1176, which corresponds to an approximately 20\% higher chance of outperforming a rival system.

Prospective studies should address the potential for biases in automated assessment frameworks and explore effective methods for mitigating such biases.

\paragraph{Data \& Training} It is important to highlight that the OASST1 dataset, used to train the {Guanaco} models, contains examples in multiple languages, and that the OA benchmark also comprises prompts in various languages. An open question for subsequent research is the extent to which this multilingual training enhances model performance on instructions presented in non-English languages and whether it accounts for the increased disparity observed between the Vicuna-13B model—exclusively trained on English data—and the {Guanaco} 33B and 65B models on the OA benchmark.

To probe for possible data leakage, given the notable performance of {Guanaco} models, we compared the OASST1 data and the Vicuna benchmark prompts. Applying fuzzy string matching techniques and manually reviewing the most similar cases, we did not detect any overlapping prompts between the datasets.

In addition, our model has only been trained using cross-entropy loss in a supervised learning setup and does not incorporate reinforcement learning from human feedback (RLHF). This presents an opportunity for future work to examine the relative advantages and limitations of cross-entropy versus RLHF-based training regimes. We anticipate that \textsc{QLoRA} facilitates such investigations at scale, obviating the requirement for prohibitive computational resources.

\section{Related Work}

\paragraph{Quantization of Large Language Models} Research in LLM quantization has been predominantly directed towards optimizing inference-time efficiency. Strategies to preserve the accuracy characteristic of 16-bit LLMs often involve specialized handling of outlier activations (as in SmoothQuant~\citep{xiao2022smoothquant} and LLM.int8()~\citep{dettmers2022llm}), with other efforts adopting advanced grouping mechanisms~\citep{park2022nuqmm, yao2022zeroquant}. Approaches utilizing lossy quantization have explored the balance between precision and regular rounding~\citep{dettmers2022case,zeng2022glm,pope2022efficiently}, as well as methods to optimize rounding choices for superior quantization accuracy~\citep{frantar2022gptq}. In terms of training dynamics, apart from this work, only SwitchBack layers~\citep{wortsman2023stable} have examined gradient-based optimization through quantized weights at model scales greater than 1B parameters.

\paragraph{Finetuning with Adapters} In our experiments, we employ Low-rank Adapters~\citep{hu2021lora} (LoRA) as our Parameter Efficient FineTuning (PEFT) method. However, there exists a wide range of alternative PEFT techniques including prompt tuning~\citep{qin2021learning,lester2021power,li2021prefix}, tuning input embeddings~\citep{an2022input}, modifying hidden states (e.g., IA$^3$)~\citep{liu2022few}, incorporating additional complete layers~\citep{houlsby2019parameter}, optimizing only bias terms~\citep{zaken2021bitfit}, applying Fisher information-guided weight masking~\citep{sung2021training}, as well as hybrid strategies~\citep{henderson2021compacter}. Our results demonstrate that LoRA adapters can achieve comparable results to standard full 16-bit finetuning. Exploration of alternative PEFT methods and their trade-offs remains an open area for future investigations.

\paragraph{Instruction Finetuning} Instruction finetuning seeks to adapt pretrained large language models (LLMs) so that they respond to prompts according to specific instructions by training on input-output pairs derived from a variety of datasets. Notable methodologies and data sources for instruction finetuning include MetaICL~\citep{min2021metaicl}, MetaTuning~\citep{zhong2021adapting}, InstructGPT~\citep{ouyang2022training}, FLAN~\citep{wei2021finetuned,chung2022scaling}, PromptSource~\citep{bach2022promptsource}, Super-NaturalInstructions~\citep{wang2022super,sanh2021multitask}, Self-instruct~\citep{wang2022self}, UnnaturalInstructions~\citep{honovich2022unnatural}, OPT-IML~\citep{iyer2022opt}, UnifiedSKG~\citep{xie2022unifiedskg}, OIG/Chip2~\citep{laion2023}, Alpaca~\citep{alpaca}, Vicuna~\citep{vicuna2023}, Koala~\citep{koala_blogpost_2023}, and Self-instruct-GPT-4~\citep{peng2023instruction}.

\paragraph{Chatbots} Instruction-tuned models are frequently deployed as conversational chatbots, where optimization may rely on Reinforcement Learning from Human Feedback (RLHF)~\citep{christiano2017deep}, or on training data synthesized using model-generated feedback (RLAIF)~\citep{bai2022constitutional}. Some of the principal models and corpora in this space are Anthropic-HH~\citep{askell2021general,bai2022training}, Open Assistant~\citep{kopf2023openassistant}, LaMDA~\citep{thoppilan2022lamda}, and Sparrow~\citep{glaese2022improving}. Although we do not implement reinforcement learning in our workflow, our top-performing system, Guanaco, has been finetuned using the Open Assistant multi-turn dialogue corpus, which was created for RLHF-based training~\citep{kopf2023openassistant}. Recent evaluation frameworks rely on GPT-4 as a substitute for labor-intensive human annotations~\citep{vicuna2023,peng2023instruction}. Our contribution refines these evaluation schemes, prioritizing reliability in assessment.

\section{Limitations and Discussion}

Our work presents evidence that \textsc{QLoRA} enables a 4-bit base model with Low-rank Adapters (LoRA) to achieve results comparable to traditional full 16-bit finetuning. Nevertheless, we have not established that \textsc{QLoRA} can match 16-bit finetuning performance at the 33B and 65B parameter scales. The significant computational costs required for experiments at these scales were beyond the scope of this study and are deferred to future research.

The evaluation protocol for instruction-finetuned models presents another constraint. While assessments were conducted on MMLU, Vicuna, and OA benchmarks, we did not include other popular evaluation suites such as BigBench, RAFT, and HELM. Consequently, it remains unclear to what extent our evaluation outcomes generalize across these additional benchmarks. On a positive note, our study undertakes an extensive analysis on MMLU and introduces novel chatbot evaluation methodologies.

The findings suggest that benchmark performance is heavily influenced by the degree of overlap between finetuning datasets and evaluation benchmarks. For instance, FLAN v2 aligns closely with MMLU but is less aligned with chatbot tasks, whereas the Chip2 dataset shows the opposite trend; the performance of respective models on MMLU and Vicuna reflects these relationships. This underscores not just the demand for improved evaluation frameworks but also the necessity of clearly defining the capabilities one seeks to measure. It raises important questions: Should the emphasis be placed on excelling in academic domains such as high school or college-level assessments, or in conversational competencies as seen in chatbot settings? Or is there an alternative target that deserves focus? The prevalence of pre-existing benchmarks often biases the research trajectory toward certain evaluation tasks, due to the relative ease of reusing benchmarks instead of developing new ones. As a community, it is vital to ensure that benchmarks accurately reflect prioritized outcomes.

In our study, although we conduct a comprehensive assessment of general chatbot abilities, we acknowledge a limitation in our responsible AI analysis of {Guanaco}. Specifically, we examine the propensity of {Guanaco}-65B to generate socially biased text, compared to other models, with the results detailed in Table~\ref{tbl:crows}. The findings show that {Guanaco}-65B attains a notably lower mean bias score compared to other pretrained counterparts, suggesting that finetuning on the OASST1 dataset can mitigate biases present in the LLaMA base model. Despite these promising indications, the question of how {Guanaco} performs in the context of other bias types remains unresolved. We leave the comprehensive exploration of bias analysis in {Guanaco} and related conversational agents to future research.

Another constraint of our work is the absence of experiments involving different quantization precisions—such as 3-bit base models—or alternative adapter strategies. Beyond LoRA, many Parameter Efficient FineTuning (PEFT) approaches have demonstrated effectiveness; however, scalability to models of larger sizes remains uncertain. While we adopted LoRA for its established stability, alternative adapters might surpass its performance. Notably, our observations suggest that finetuning after quantization can recover a substantial portion of information lost in quantization, potentially permitting more aggressive approaches. For instance, applying 3-bit GPTQ quantization to the base model in conjunction with LoRA finetuning may deliver performance levels on par with full 16-bit finetuning after adaptation.

\section{Broader Impacts}

Our \textsc{QLoRA} finetuning technique stands out as the inaugural method capable of finetuning models with 33B parameters using a single consumer-grade GPU, and models with 65B parameters using a single professional GPU, all while maintaining performance comparable to traditional full finetuning. We have shown that our top-performing 33B model, trained on the Open Assistant dataset, achieves competitive results with ChatGPT on the Vicuna benchmark. As instruction finetuning remains essential for converting generic pretrained LLMs into effective ChatGPT-like conversational agents, our approach substantially democratizes the technology, particularly benefiting resource-limited researchers and advancing the accessibility of advanced NLP systems. In this light, \textsc{QLoRA} acts as a leveling mechanism, reducing the disparity between large industrial labs and smaller research teams relying on consumer hardware.

Another avenue for significant impact is the adaptation of these models for mobile devices. We propose that our \textsc{QLoRA} technique may represent a pivotal advancement, permitting the finetuning of LLMs directly on smartphones and in other environments with constrained computational resources. Although previous research has demonstrated that 7B parameter models can be executed on mobile phones, \textsc{QLoRA} is, to our knowledge, the inaugural approach capable of supporting the finetuning process for these models on such devices. Our estimates suggest that with a device like the iPhone 12 Plus, it is feasible to finetune on approximately 3 million tokens overnight while the device charges. Despite the fact that these finetuned 7B models do not yet match the performance levels of systems such as ChatGPT, their effectiveness appears sufficient to enable innovative applications previously hindered by concerns regarding user privacy or the limitations of LLM capabilities. The use of \textsc{QLoRA} paves the way for privacy-respecting applications of LLMs, granting users greater control over their data and models while also reducing barriers to LLM deployment.

Nevertheless, it is important to recognize that the ability to finetune LLMs can be misused, given its dual-use nature. Although broad deployment of LLMs entails recognized risks \citep{bommasani2021opportunities,bender2021dangers}, we maintain that democratizing access to this increasingly pervasive technology will foster more robust and independent oversight, as opposed to concentrating the control of LLMs within major corporations that may limit transparency by withholding models and source code.

In summary, we are optimistic that \textsc{QLoRA} will substantially advance the accessibility and practicality of high-quality LLM finetuning, thus expanding its positive influence.